\documentclass[10pt]{article}

\usepackage[utf8]{inputenc}

\usepackage[T1]{fontenc}

\usepackage{amsmath}

\usepackage{amsfonts}

\usepackage{amssymb}

\usepackage[version=4]{mhchem}

\usepackage{stmaryrd}



\begin{document}

Славнова - Тейлора тождества, связывающие различные функции Грина и являющиеся следствием калибровочной инвариантности лагранжиана (2). Перенормированные квантовохромодинамич. функции Грина зависят от кваркглюонной константы связи $\alpha_{s}=g^{2} / 4 \pi$, масс кварков и точки нормировки $\mu$. Зависимость от последней возникает вследствие известной конечной неопределенности в выборе перенормировочных констант (см. [5]).



Примечательной особенностью К. х. является свойство асимптотической свободы: эффективная константа связи $\bar{\alpha}_{s}(\mu)$, подчиняющаяся стандартному уравнению ренормализационной группы



\[

\mu^{2} \frac{d}{d \mu}=\bar{\alpha}_{s}(\mu)=\beta\left(\alpha_{s}\right)

\]



убывает в ультрафиолетовой области. Из явных вычислений функции $\beta\left(\bar{\alpha}_{S}\right)$ в двух первых порядках теории возмущений следует, что



\[

\bar{\alpha}_{s}(\mu)=\frac{12 \pi}{\left(33-2 n_{f}\right) \ln \left(\mu^{2} / \Lambda^{2}\right)}\left(1-\frac{6\left(153-19 n_{f}\right) \ln \left(\ln \left(\mu^{2} / \Lambda^{2}\right)\right)}{\left(33-2 n_{f}\right) \ln ^{2}\left(\mu^{2} / \Lambda^{2}\right)}\right),

\]



где $n_{f}$ - число кварков с массой менышей $\mu, \Lambda-$ так наз. п а ра м е т р шкалы К. х. Физич. смысл $\Lambda$ состоит в том, что числовое значение этого параметра характеризует масштаб области сильной связи в К. х. При выводе формулы (6) полностью игнорировались эффекты масс кварков, то есть рассматривалась идеализированная ситуация, когда влиянием кварков с массой больше $\mu$ полностью пренебрегалось. В этом приближении зависящие от $n_{f}$ члены в формуле (6) меняются дискретным образом, когда пересекаются кварковые пороги. В области болыших импульсных передач $\mu$ эффективная константа связи $\bar{\alpha}_{s}(Q)=\bar{\alpha}_{3}(\mu=0)$ логарифмически падает, то есть в этом пределе К. х. превращается в теорию слабо взаимодействующих между собой кварков и глюонов. Этот факт позволяет использовать улучшенную с помощью аппарата ренормализационной группы (см. [6]) теорию возмущений для надежного нахождения функций Грина в области больших импульсов $Q$ (то есть когда $\alpha_{s}(Q)<c$ ), что соответствует на координатном языке малым расстояниям. В отличие от квантовой электродинамики, сама возможность осмысленного применения теории возмущений в К. х. базируется на предварительном использовании соотношений ренормализационной группы. Свойство асимптотич. свободы позволило физически обосновать справедливость партонной картины сильных взаимодействий, в рамках к-рой адроны при определенных условиях можно рассматривать как совокупности почти не взаимодействующих между собой составляющих - партонов.



Как следует из формулы (6), при уменьшении $Q$ эффективная константа связи $\bar{\alpha}_{s}(Q)$ растет и $\geqslant 1$ при $Q$, приближаюшемся к $\Lambda$. Точное поведение $\bar{\alpha}_{S}(Q)$ при $Q \leqslant \Lambda$ в настоящее время неизвестно; сама теория возмущений и, следовательно, формула (6) неприменимы в этой области вследствие немалости эффективной константы связи. Во всяком случае, это означает, что в К. х. взаимодействие при малых переданных импульсах, то есть на болыших расстояниях, в отличие от квантовой электродинамики, становится сильным, что не позволяет непосредственно использовать лагранжиан (4) для определения наблюдаемого спектра одночастичных состояний в этой теории. Согласно гипотезе конфайнмента, все физич. состояния в К. х. бесцветны, то есть являются синглетами по отношению к цветовой группе $S U_{c}$ (3). Кварки и глюоны, реализуя фундаментальное и, соответственно, присоединенное представления группы $S U_{c}(3)$, не могут существовать в свободном состоянии как физич. частицы. При этом экспериментально наблюдаемые сильно взаимодействующие частицы - адроны - представляют собой бесцветные связанные состояния кварков и глюонов. Одним из следствий гипотезы конфайнмента является естественное объяснение отсутствия дробно заряженных адронов, что подтверждается всей существующей совокупностью экспериментальных данных. В настоящее время строгое доказательство этого явления отсутствует и детальный механизм удержания кварков внутри адронов неясен. Одно из популярных качественных объяснений гипотезы конфайнмента состоит в том, что на больших по сравнению с $\hbar c / \Lambda$ расстояниях сила притяжения кварка и антикварка друг к другу не зависит от расстояния между ними и, следовательно, потенциальная энергия их взаимодействия пропорциональна этому расстоянию. Именно такое поведение потенциальной энергии было найдено в рамках так наз. решеточного приближения в К. х. (см. Вильсона критерий конфайнмента).



Опосредованный характер К. х. и трудности с доказательством гипотезы конфайнмента значительно усложняют задачу экспериментальной проверки достоверности этой теории и нахождения значения фундаментального параметра $\Lambda$. Несмотря на это, к настоящему времени имеется значительное число количественных предсказаний К. х., подтвержденных экспериментально. Большая часть этих результатов получена с помощью теории возмущений. Из предыдущего следует, что теория возмущений К. х. могла бы быть применима для количественного описания так наз. жестких процессов с большим переданным импульсом $Q$, для к-рых эффективная константа связи $\bar{\alpha}_{s}(Q)$ достаточно мала. Естественно ожидать, что в итоге должно быть возможным воспроизведение в случае жестких процессов основных результатов кварк-партонной модели и найти к ним поправки, связанные с высшими членами теории возмушений по $\bar{\alpha}_{s}$. Главная, до сих пор не решенная окончательно трудность в реализации этой программы заключается в том, что информацию о реальных адронных процессах приходится извлекать из функций Грина кварковых и глюонных полей (возможно, со вставками локальных операторов, составленных из этих полей), рассматриваемых в области болыших евклидовых импульсных передач, к-рые лишь и могут надежно вычисляться в рамках теории возмущений. Тем не менее, в нек-рых важных частных случаях, к к-рым относятся, в частности, полное сечение электрон-позитронной аннигиляции в адроны и инклюзивные глубоко-неупругие лептон-адронные реакции, эту проблему можно обойти. В результате оказалось возможным не только обосновать наблюдаемое экспериментально приблизительно автомодельное поведение соответствующих сечений, но и предсказать позднее обнаруженные экспериментально логарифмич. отклонения от этого поведения, возникающие вследствие вкладов высших порядков теории возмущений.



K сожалению, метод операторных разложений непосредственно не применим к рассмотрению более сложных жестких процессов, таких, как рождение $\mu^{+} \mu^{-}$пар, упругих электромагнитных адронных формфакторов или рождения частиц с большим поперечным импульсом в адронных столкновениях, закономерности асимптотич. поведения к-рых хорошо описываются правилами кваркового счета (см. [7]). В этом случае вместо разложений Вильсона используют прямой анализ диаграмм Фейнмана, даюших вклад в изучаемый процесс. Несмотря на то что такой подход сушественно более трудоемок и сталкивается с определенными технич. трудностями, с его помощью также удается подтвердить предсказания кварк-партонной модели и формул кваркового счета и вычислить по крайней мере первый неисчезающий поправочный член порядка $\bar{\alpha}_{s}$ для жестких процессов. Во всех случаях, в к-рых предсказания К. х. для жестких процессов сравнивались с экспериментальными данными, наблюдалось взаимное согласие. При этом определяемое в разных процессах значение масштабного параметра $\Lambda$ оказывается лежащим в относительно узком интервале: $\Lambda=100 \div 300$ МэВ (если принять во внима-





\end{document}